\subsection{nml\_debug}
\label{SEC:nml_debug}
\inp*[nmlDebug]{}[nml\_debug][]

This \hyperref[SEC:User_Input]{namelist block}
is used by developers for debugging.

\begin{nmlblocklst}[
,label={LST:nml_debug}
,caption={Namelist nml\_debug.}
]
&nml_debug
  Output_Limiter_Statistics          = F
  Output_Filter_Diagonals            = F
  Output_Plot3d_To_Tecplot           = F
  Output_Bface_Array                 = F
  Dump_Memory_Usage                  = F
  Check_Flux_Point_Coordinates       = F
  Force_Nd_Subroutines               = F
  Newton_Convergence_Goal            = 1.0e-12
  Chemistry_Sub_Iterations           = -1
  Traceback_Unconverged_Newton_Temps = F
  Ignore_Thermo_Newton_Warnings      = F
/
\end{nmlblocklst}

%==========================================================================

\begin{nmldescitem}
Output_Limiter_Statistics = F
\end{nmldescitem}
\inp*[nmlDebug::OutputLimiterStatistics]{}[Output\_Limiter\_Statistics][]

\begin{nmldescbody}
This outputs the number of limited\slash projected cells each time step when
\inp[nmlLimiter::LocalLimiterOption]{Local\_Limiter\_Option} is set to use
projection.
\end{nmldescbody}

%==========================================================================

\begin{nmldescitem}
Output_Filter_Diagonals = F
\end{nmldescitem}
\inp*[nmlDebug::OutputFilterDiagonals]{}[Output\_Filter\_Diagonals][]

\begin{nmldescbody}
Specifies whether to output text files containing the filter diagonals and
filter matrices for all geometry types after they are generated in the
preprocessor.
%
This is used to inspect values for possible errors in the filter
definition.
\end{nmldescbody}

%==========================================================================

\begin{nmldescitem}
Output_Plot3d_To_Tecplot = F
\end{nmldescitem}
\inp*[nmlDebug::OutputPlot3dToTecplot]{}[Output\_Plot3d\_To\_Tecplot][]

\begin{nmldescbody}
Specifies whether to output a Tecplot file of the input Plot3D grid file
after reading it.
%
This can be used to inspect whether the grid file was read correctly.
\end{nmldescbody}

%==========================================================================

\begin{nmldescitem}
Output_Bface_Array = F
\end{nmldescitem}
\inp*[nmlDebug::OutputBfaceArray]{}[Output\_Bface\_Array][]

\begin{nmldescbody}
Specifies whether to have each MPI process output the contents of its
own local bface array to the text file ``bface.[MPI process \#].dat''.
%
The bface array contains the boundary conditions and connectivity
information for each boundary face.
\end{nmldescbody}

%==========================================================================

\begin{nmldescitem}
Dump_Memory_Usage = F
\end{nmldescitem}
\inp*[nmlDebug::DumpMemoryUsage]{}[Dump\_Memory\_Usage][]

\begin{nmldescbody}
Specifies whether to profile and output the memory usage of GFR.
\end{nmldescbody}

%==========================================================================

\begin{nmldescitem}
Check_Flux_Point_Coordinates = F
\end{nmldescitem}
\inp*[nmlDebug::CheckFluxPointCoordinates]{}[Check\_Flux\_Point\_Coordinates][]

\begin{nmldescbody}
Flag used to output files to check that co-located face quantities
are consistent in terms of interpolation between two cells sharing a face.
\end{nmldescbody}

%==========================================================================

\FINALFIX{\newpage}
\begin{nmldescitem}
Force_Nd_Subroutines = F
\end{nmldescitem}
\inp*[nmlDebug::ForceNdSubroutines]{}[Force\_Nd\_Subroutines][]

\begin{nmldescbody}
GFR normally uses optimized 1D, 2D, and 3D versions of various
subroutines.
%
This input forces the use of the generalized ``nD'' versions of the
various subroutines that loop over the number of equations instead
of hardcoding in each equation.
\end{nmldescbody}

%==========================================================================

\begin{nmldescitem}
Newton_Convergence_Goal = 1.0e-12
\end{nmldescitem}
\inp*[nmlDebug::NewtonConvergenceGoal]{}[Newton\_Convergence\_Goal][]

\begin{nmldescbody}
This specifies the convergence goal for all of the Newton iterations on
temperature.
\end{nmldescbody}

%==========================================================================

\begin{nmldescitem}
Chemistry_Sub_Iterations = -1
\end{nmldescitem}
\inp*[nmlDebug::ChemistrySubIterations]{}[Chemistry\_Sub\_Iterations][]

\begin{nmldescbody}
Number of chemistry sub-iterations to use for the KIVA scheme source terms.
%
\NOTE
The number of sub-iterations is set to 10 if
$\inptext{Chemistry\_Sub\_Iterations}<1$.
\end{nmldescbody}

%==========================================================================

\begin{nmldescitem}
Traceback_Unconverged_Newton_Temps = F
\end{nmldescitem}
\inp*[nmlDebug::TracebackUnconvergedNewtonTemps]{}[Traceback\_Unconverged\_Newton\_Temps][]

\begin{nmldescbody}
If running a multi-component gas mixture, setting this input to true
will write a traceback to stdout whenever the Newton solve for
temperature fails to meet the convergence goal.

\NOTE
This only works if using the Intel Fortran compiler and
the namelist input variable
\inp[nmlDebug::IgnoreThermoNewtonWarnings]{Ignore\_Thermo\_Newton\_Warnings}
is set to false.
\end{nmldescbody}

%==========================================================================

\begin{nmldescitem}
Ignore_Thermo_Newton_Warnings = F
\end{nmldescitem}
\inp*[nmlDebug::IgnoreThermoNewtonWarnings]{}[Ignore\_Thermo\_Newton\_Warnings][]

\begin{nmldescbody}
If running a multi-component gas mixture, setting this to true will prevent
any warnings from being written to stdout whenever the Newton solve for
temperature fails to converge.
\end{nmldescbody}

%==========================================================================

